% V19 recollation: French PDF physical 50 (last paragraph)--52; English 34 (last paragraph)--36 (before section 9).
% Replaces only DOT_IMPORT 11 and 12 in v18_first_8_14_onward.tex.
% Confirmed corrections and additional mathematical hypotheses are distinguished in footnotes.
\begingroup
\subsection*{8.18之后的线性系讨论}

还须把8.18的条件(a)、(b)写得明确一些。这使我们一般地考察如下情形。
假定 $X$ 在 $k$ 上几何不可约，并给定 $P^\vee$ 的一个线性子簇 $L$。
这对应于研究经过给定点 $x$，或者在给定光滑点与 $X$ 相切的超平面截面。
因此，$L$ 由那些包含 $P$ 中某个线性子簇 $L^0$ 的超平面组成；
在上述两种情形中，$L^0$ 分别是一个点，或者 $X$ 在一个光滑点的切空间的像。
我们要问：对于 $L$ 的泛点 $\xi$，因而对于 $L$ 的某个非空开子集的所有点，
$Y_\xi$ 是否几何不可约，且维数为 $\dim X-1$？\nobreak\cmt{M19-8E01}{本段据法文物理50页末至51页上半重译。法稿明确写8.18，英编稿34页误作1.18；
法51的参数空间为 $P^\vee$，中心为 $L^0$，英编稿35页分别漏了对偶标记和上标0。
法50把“$X$ 整、$k$ 代数闭”划改为“$X$ 在 $k$ 上几何不可约”；法51把旧词
\emph{simple} 改为 \emph{lisse}，并把截面条件中原“几何整”的字样划掉，改为“几何不可约”。
正文采用可辨改写层。英编稿所留 \emph{de ravoi} 只是法文 \emph{de savoir} 的误读，
这里完整译为“要问……是否”。}

这是 Bertini 定理的一个变式，应当列在第3节中，并且完全可以用同一方法处理
（或者，愿意的话，也可以归结到那里）。

保证泛截面不包含整个 $X$ 的条件是
\[
 f(X)\not\subset L^0,
\]
也就是说，
\[
 X'=f^{-1}(P\setminus L^0)
\]
是 $X$ 的稠密开子集。若泛截面非空，其维数便是 $\dim X-1$。\nobreak\cmt{M19-8E02}{法51直接说“维数问题仅表达为 $f(X)\not\subset L^0$”，英编稿又把包含符号误排。
该条件保证泛超平面所给的方程在 $X_{\mathrm{red}}$ 上不恒为零，却不单独保证泛截面非空；
例如 $f(X)$ 为中心之外的一个点时，泛截面为空。正文把非空性与维数结论分开；
后面的正维投影像条件保证所需的非空性。}

令 $Q$ 为参数空间 $L$ 的对偶射影空间。注意，若定义 $L^0$ 的线性形式组成
$E$ 的子空间 $F^0$，亦即 $L^0=\mathbf P(E/F^0)$，则
\[
 Q=\mathbf P(F^0).
\]
考虑由 $F^0\hookrightarrow E$ 导出的典范态射（参看第二章）
\[
 u:P\setminus L^0\longrightarrow Q,
\]
并考虑
\[
 g=u f':\ f^{-1}(P\setminus L^0)=X'\longrightarrow Q.
\]
于是 $L\simeq Q^\vee$\cmt{M19-8E03}{法51首次命名 $Q$ 的字母有覆打，其措辞若依英编稿直读为“$Q$ 是经过 $L^0$
的超平面的参数空间”，便同后面的 $L\simeq Q^\vee$ 混淆。
正文按法稿明确保留的 $Q=\mathbf P(F^0)$、$F^0\to E$ 和 $L\simeq Q^\vee$
统一约定，并明示本篇采用的可逆商约定；不是声称覆打处已完全辨清。
法稿在 $\mathbf P(F^0)$ 的括号内划掉 $E/$，又把导出映射中的旧商记号划去，改为 $F^0\to E$。
英编稿的 $F^-\to E$ 以及 $f^{-1}(P-L)$ 均据法稿订正。
这里 $f'$ 是 $f$ 在 $X'$ 上的限制，$X'_\xi=Y_\xi\cap X'_{k(\xi)}$；
原稿在本段用 $X_\xi$ 表示先前的 $Y_\xi$。}，而 $X'_\xi$（$\xi\in L$）这一族，正是相对于态射 $g$
的超平面截面族。

另一方面，要从 $X'_\xi$ 的不可约性得到整个 $Y_\xi$ 的不可约性，还须保证前者在后者中稠密。
在这里，一个足够的条件是基点集 $B=f^{-1}(L^0)$ 没有余维为1的不可约分支，即
\[
 \dim B\leq \dim X-2
\]
（空集也允许）。在这个附加条件下，对充分一般的 $\xi\in L$，非空的 $X'_\xi$
在 $Y_\xi$ 中稠密；这一点在扩张到底域的代数闭包后仍成立。因此，$X'_\xi$ 几何不可约，
当且仅当 $Y_\xi$ 几何不可约。\nobreak\cmt{M19-8E04}{这是对原论证补明的译校条件，不是法稿中恢复出的缺字。
法51无条件声称：对一般 $\xi$，$X'_\xi$ 在 $X_\xi$ 中稠密，因而二者的几何不可约性等价；
英编稿35页又漏去了这些截面的下标。仅有 $X'$ 在 $X$ 中稠密不足以推出此断言。
反例为二次 Veronese 浸入
$[u:v:w]\mapsto[u^2:uv:uw:v^2:vw:w^2]$，取中心 $L^0=V(z_0,z_1,z_2)$。
泛截面为 $u(au+bv+cw)=0$，含固定直线 $u=0$，却有 $\dim g(X')=2$。
正文的附加条件排除了这个问题；完整的维数论证和反例见核查报告 M19-8E。
若只研究去掉基点集后的截面 $X'_\xi$，无需此附加条件。}

在此前提下，Bertini--Zariski 定理表明，只要
\[
 \dim g(X')\geq2,
\]
就得到所希望的不可约性结论。

（说实在，顺便还可以给3.1补充一个逆命题：设 $X$ 几何不可约，并把 $Y$ 理解为泛超平面截面。
在几何基变换后，它几何不可约，当且仅当以下两种情形之一成立：
或者 $\dim f(X)\geq2$；或者 $\dim f(X)=1$，且像包含在一条直线 $D$ 中，
而到 $D$ 的态射的泛纤维几何不可约。这里直线及后一泛纤维条件在
$\overline k$ 上理解。）\nobreak\cmt{M19-8E05}{法51的第一维数分支处有改动，现存符号形似“$\ne2$”，不能据该字形确认终稿；
英编稿35页写成“$=2$”。正文采用 $\geq2$，依据是3.1的适用范围及此处的几何判据；
例如 $f=\mathrm{id}_{\mathbf P^3}$ 的泛截面是几何不可约的射影平面，不能被“$=2$”排除。
法稿直线旁明写“定义在 $\overline k$ 上”，英编稿漏横线；泛纤维条件则由法稿页脚手写补入。
正文把相关对象同时基变换到 $\overline k$，避免写出从未作基变换的 $X$ 到只在
$\overline k$ 上定义的 $D$ 的态射。像为曲线时的判据证明另列报告，属于译校补证；
作者在这里仅提出补写逆命题，没有给出这份证明。}

在上述稠密性条件下，把这一判据用于 $g$，也就得到当前带有 $L$ 的版本中，
泛点 $\xi\in L$ 所给的 $Y_\xi$ 几何不可约的充要条件。

从“cunutesque”〔原词含义待考〕的观点出发，用域论的语言，可以如下把条件写成超越次数的条件。\nobreak\cmt{M19-8E06}{法52可读到这个非常规词，英编稿亦以问号保留；不以“计算的”或其他猜测替代。
紧接“条件”之后的 \emph{fondamentale}（基本的）已划掉，故不进入正文。}
选择一个既不包含 $L^0$、也不包含 $f(X)$ 的“无穷远超平面”，然后在其补集中讨论；
也就是说，实质上在标准仿射开集上讨论。\nobreak\cmt{M19-8E07}{法稿沿用“不包含 $X$”这一措辞；一般 $f:X\to P$ 下，它指不包含整个像 $f(X)$，
并在该仿射开集的逆像上考察 $X$。不是断言整个 $X$ 已包含于一个仿射开集，也不是新增 $f$ 为浸入的假设。}
取在 $L^0$ 上消失的线性形式空间的一组基
\[
 T_1,\ldots,T_p,\qquad p=\operatorname{codim}(L^0,P),
\]
并考虑它们在 $X$ 的分式域 $K$ 中的拉回 $t_1,\ldots,t_p$（这里假定 $X$ 整）。
至少有一个 $t_i$ 非零，不妨设 $t_1\ne0$。于是考虑
\[
 a_1=t_2/t_1,\quad\ldots,\quad a_{p-1}=t_p/t_1.
\]
那么 $\dim g(X')$ 正是子域
\[
 k(a_1,\ldots,a_{p-1})\subset K
\]
对 $k$ 的超越次数。\nobreak\cmt{M19-8E08}{法52明确写中心 $L^0$ 和小写 $k(a_1,\ldots,a_{p-1})$，据此订正英编稿35页的
“在 $L$ 上消失”以及 $K(a_1,\ldots,a_{p-1})$。在所选仿射开集上，齐次线性形式经平凡化成为函数。
这些坐标比生成投影像闭包的函数域，因此这也是维数等式的直接依据；像的维数均指其闭包的维数。}

因此，若这个超越次数至少为2，就满足前述条件。若它为1，则应要求：在 $\overline k$ 上，
$f(X)$ 包含在一个含 $L^0$、维数比 $L^0$ 大\emph{两}维的线性子空间中，
并且 $g:X'\to\overline{g(X')}$ 的泛纤维几何不可约。\nobreak\cmt{M19-8E09}{法52写的是“包含 $L^0$，而维数多一维”
（\emph{de dimension une de plus}）；英编稿35--36页进一步误译为“维数至多1”。
两者都不能给出这里的条件。若 $q=\dim L^0$，投影 $u$ 的纤维为 $q+1$ 维仿射空间；
一条直线 $D\subset Q$ 的逆像的射影闭包是含 $L^0$ 的 $q+2$ 维线性子空间。
正文因此订为多两维。最小例子是 $X=P=\mathbf P^2$、$f$ 为恒等态射、$L^0$ 为一点：
投影像为直线，但 $f(X)$ 不是含该点的一条直线。法稿此处的底域明确为 $\overline k$，
并在行间手写加入到像闭包的泛纤维条件；英编稿漏去底域横线。}

设 $L^0$ 的维数为 $q$。那么
\[
 u:P\setminus L^0\longrightarrow Q
\]
的纤维维数是 $q+1$。它限制到 $\overline{f(X)}\setminus L^0$ 后，各纤维的维数至多为
$q+1$，因而
\[
 \dim g(X')\geq\dim f(X)-(q+1).
\]
所以，只要
\[
 \dim f(X)\geq q+3,
\]
就满足关于 $g(X')$ 的维数条件。\nobreak\cmt{M19-8E10}{法52行间手写的原句是“所以 $g$ 的纤维维数 $\leq q+1$”，英编稿36页亦如此。
若 $f$ 本身具有正维纤维，这句话不成立；例如 $f:\mathbf P^N\times\mathbf P^m\to\mathbf P^N$
为投影时，$g$ 的纤维维数是 $q+1+m$。所需维数不等式应对像闭包上的投影来证明，
正文据此订正，保留原有结论。法稿记中心维数为 $q$，英编稿此处的 $g$ 是误字。}

在基点集没有余维1分支这一附加条件下，$q=0$ 时就得到8.17(a)所指的事实。\nobreak\cmt{M19-8E11}{法52直接说“$q=0$ 时，便得到8.17(a)指出的事实”，未附限制。
如果仍只假设 $X$ 几何不可约且 $\dim f(X)\geq3$，该说法过强：
$\operatorname{Bl}_p\mathbf P^3\to\mathbf P^3$ 中，所有经过 $p$ 的超平面的总变换都含例外除子，
并含不同的严格变换，因而可约；这些参数本身构成一个超平面。
这与8.18另作的非分歧性假设是不同的范围。正文明确保留本段所需的基点条件，不据此重新断言8.17(a)的宽泛版本已经证明。}

回到原8.18(a)的不可约性要求，作者在这里给出如下提示：对于 $x\in T$，
只要 $X$ 相对于 $f$“在 $x$ 处不是锥形的”，便满足该条件；这里“锥形”的意义应当是明显的。\nobreak\cmt{M19-8E12}{法52明确是大写 $X$“在 $x$ 处不是锥形的”，英编稿误成小写 $x$。
“意义明显”是作者的表述；此处未定义该相对概念，也没有写出验证。
正文保留其提示身份，不把它视为已证明的形式判据，尤其不能直接用来验证本版8.18中另经加强的几何整条件。}

还须考虑，是否把这些最后的 Bertini 式讨论放到下一节“更换射影浸入”中更合适。
\par
\endgroup
